\begin{proof}
\begin{enumerate}[label=\arabic*.]
\item By \cref{maps:src:F.2.9,maps:src:F.2.15}, the cohomology
system $V$ has a one-dimensional invariant space $I$, no invariant
covectors, and two-dimensional finite fixed spaces $F_1,F_2$.
The finite orders are among the coprime pairs of
\cref{maps:src:F.2.8}.  At a finite quotient, averaging identifies
the degree-two specialization with the fixed part of $\Lambda^2V$.
At the cusp, \cref{night:P1.11} gives the actual surjective homology
collapse with kernel the saturation of
$(\Lambda^2M_0-I)\Lambda^2\Lambda$.  Its dual has image
$\ker(\Lambda^2M_0^t-I)$ under the evaluation pairing.  Consequently
the natural specialization morphism identifies
\[
 R^2g_*\mathbb Q=j_*\Lambda^2V.
\]
The term $H^0(j_*\Lambda^2V)$ has no incoming Leray differential,
and its only possible outgoing differential targets
$H^2(j_*V)=H^0(V^*)^*=0$.  It survives to $H^2(M;\mathbb Q)=0$.
Thus $a_2':=\dim(\Lambda^2V)^\Gamma=0$.

\item Let $\beta:V\to V^*$ be an invariant alternating form.
Since $\beta(I)\subset(V^*)^\Gamma=0$, its radical contains $I$.
For the finite averaging projection $P_i:V\to F_i$, invariance gives
$\beta(u,x)=\beta(u,P_ix)$ when $u\in F_i$.
The alternating restriction to the two-dimensional space $F_i$
has $I$ in its radical and is therefore zero.  Hence
$F_1+F_2\subset\operatorname{rad}\beta$.  This sum has dimension
three because $F_1\cap F_2=I$.  The alternating rank is even,
so every such form has rank at most two.

Suppose $\beta,\gamma$ were independent invariant forms.
They and their sum have rank two, so
$\beta\wedge\gamma=0$.  Their actual image planes in $V^*$
are distinct invariant two-planes meeting in a rational line.
Each finite generator acts on that line by a sign.  At least one
finite order is odd, and the unipotent product acts trivially on the
line; both signs are therefore $+1$.  This contradicts
$(V^*)^\Gamma=0$.  It follows that
$b_2':=\dim(\Lambda^2V^*)^\Gamma\leq1$.

\item At the cusp, the exterior square of three size-two unipotent
blocks has fixed rank nine: three fixed vectors come from their
individual exterior squares and two from each of the three pairwise
tensor products.  If $d_i$ denotes a finite fixed rank on
$\Lambda^2V$, the invariant-sheaf Euler identity gives
\[
 -15+d_1+d_2+9
 =a_2'-h^1(j_*\Lambda^2V)+b_2'\leq1.
 \tag{F.40}\label{maps:eq:F.40}
\]
Here the degree-two group is dual to
$(\Lambda^2V^*)^\Gamma$, by compact-support duality and the
extension-by-zero sequence.

\item The two-dimensional unit eigenspace leaves four dimensions
for the nonunit cyclotomic factors.  Local complex compatibility
requires even multiplicity of $-1$.  Partitioning four into the
cyclotomic degrees gives the following complete list:
\[
\begin{array}{c|l}
3&\Phi_3^2\\
4&\Phi_4^2,\ \Phi_2^2\Phi_4\\
5&\Phi_5\\
6&\Phi_6^2,\ \Phi_2^2\Phi_3,\ \Phi_2^2\Phi_6,\ \Phi_3\Phi_6\\
8&\Phi_8\\
10&\Phi_{10}\\
12&\Phi_{12},\ \Phi_3\Phi_4,\ \Phi_4\Phi_6.
\end{array}
\]
The unit eigenspace contributes one fixed bivector.  A repeated
nonreal quadratic pair contributes four further fixed bivectors,
so $d_i=5$ for $\Phi_3^2,\Phi_4^2,\Phi_6^2$.
Every other listed factor contributes two, so $d_i=3$.
Inequality \eqref{maps:eq:F.40} excludes every occurrence of
$d_i=5$.  In particular all order-three pairs disappear.
The coprime order list of \cref{maps:src:F.2.8} therefore reduces
to $(5,4),(5,6),(5,8),(5,12)$; their remaining factors are exactly
the eight entries of \eqref{maps:eq:F.39}.
Each pair contains the order-five factor, so
\cref{maps:src:F.2.9} proves irreducibility of $V/I$.
\end{enumerate}
\end{proof}
